\section{Лекция 6. Обобщенно-однородные, линейные дифференциальные уравнения, уравнения Бернулли и Риккати}
\noeqref{eq:l6_gen_hom_sub, eq:l6_gen_hom_ex, eq:l6_lnde, eq:l6_lode}
\noeqref{eq:l6_lode_sol, eq:l6_lagr_sub, eq:l6_lnde_struct, eq:l6_bern_sub}
\noeqref{eq:l6_bern_eq, eq:l6_bern_z, eq:l6_bern_lin, eq:l6_riccati}
\noeqref{eq:l6_riccati_spec, eq:l6_riccati_sub, eq:l6_riccati_bern}
\noeqref{eq:l6_riccati_u, eq:l6_riccati_lin, eq:l6_task1_lin, eq:l6_task1_u1}
\noeqref{eq:l6_task1_v, eq:l6_task2_hom}

\begin{flushright}
    12.10.26
\end{flushright}

\subsection{Обобщенно-однородные дифференциальные уравнения}

Некоторые нелинейные уравнения, не являющиеся однородными в смысле
определения 5.2 из Лекции 5, сводятся к однородным с помощью степенной замены
искомой функции.

\begin{definition}
    Дифференциальное уравнение 1-го порядка называется
    \textdef{обобщенно-однородным}, если существует такое вещественное число
    $m \in \R$, что при приписывании переменной $x$ измерения (веса) $1$,
    функции $y$ — измерения $m$, а производной $y' = \frac{dy}{dx}$ —
    измерения $m - 1$, все слагаемые уравнения приобретают одинаковое суммарное
    измерение.
\end{definition}

\textbf{Алгоритм решения:}
Выполним степенную подстановку:
\begin{equation} \label{eq:l6_gen_hom_sub}
    y = z^m \implies y' = m z^{m-1} z'. \tag{1}
\end{equation}
Приравнивая суммарные показатели степени (при которых $x$ и $z$ имеют вес $1$)
у всех слагаемых уравнения, находим показатель $m$. После подстановки
найденного $m$ и сокращения общих множителей получается классическое однородное
уравнение вида $z' = f(z/x)$ из Лекции 5.

\begin{example} ~
    \begin{equation} \label{eq:l6_gen_hom_ex}
        2 x^2 y' = y^3 + x y. \tag{2}
    \end{equation}

    \begin{proof}[Решение]~\\
        \begin{quote}
            \textbf{Шаг 1. Определение показателя $m$.}\\
            Подставляя $y = z^m$ и $y' = m z^{m-1} z'$, получаем:
            \[
                2 m x^2 z^{m-1} z' = z^{3m} + x z^m.
            \]
            Выпишем суммарные измерения слагаемых (считая веса $x$ и $z$
            равными $1$, а вес $z'$ равным $0$):
            \begin{itemize}
                \item для $2 m x^2 z^{m-1} z'$: измерение равно $2 + (m - 1) = m + 1$;
                \item для $z^{3m}$: измерение равно $3m$;
                \item для $x z^m$: измерение равно $1 + m$.
            \end{itemize}
            Приравнивая измерения, находим:
            \[
                m + 1 = 3m \implies 2m = 1 \implies m = \frac{1}{2}.
            \]

            \textbf{Шаг 2. Приведение к однородному уравнению.}\\
            Полагаем $y = z^{1/2} = \sqrt{z}$ ($z > 0$), тогда
            $y' = \frac{z'}{2\sqrt{z}}$:
            \[
                2 x^2 \frac{z'}{2\sqrt{z}} = z^{3/2} + x \sqrt{z}
                \implies x^2 z' = z^2 + x z \implies
                z' = \left(\frac{z}{x}\right)^2 + \frac{z}{x}.
            \]

            \textbf{Шаг 3. Решение однородного уравнения и обратная замена.}\\
            Замена $z = u x \implies z' = u' x + u$ даёт:
            \[
                u' x + u = u^2 + u \implies \frac{du}{u^2} = \frac{dx}{x}
                \implies -\frac{1}{u} = \ln \abs{x} + C.
            \]
            Возвращаясь к $z = u x$ и $y^2 = z$, получаем:
            \[
                z = \frac{-x}{\ln \abs{x} + C} \implies
                y(x) = \pm \sqrt{\frac{-x}{\ln \abs{x} + C}}.
            \]
            Потерянный при делении на $u^2$ корень $u = 0$ даёт решение
            $y(x) \equiv 0$, которое выписывается отдельно.

            \textbf{Ответ:} $y = \pm \sqrt{\dfrac{-x}{\ln \abs{x} + C}}$, а
            также $y = 0$.
        \end{quote}
    \end{proof}
\end{example}

\subsection{Линейные дифференциальные уравнения 1-го порядка}

\begin{definition}
    \textdef{Линейным неоднородным дифференциальным уравнением (ЛНДУ)} 1-го
    порядка называется уравнение, линейное относительно искомой функции $y$ и
    её производной $y'$:
    \begin{equation} \label{eq:l6_lnde}
        a_0(x) y' + a_1(x) y = b(x), \tag{3}
    \end{equation}
    где коэффициенты и правая часть непрерывны на интервале $I \subset \R$:
    $a_0, a_1, b \in C(I)$, причём $a_0(x) \neq 0$ для всех $x \in I$.
\end{definition}

\begin{definition}
    Если $b(x) \equiv 0$ на $I$, то уравнение называется \textdef{линейным
    однородным дифференциальным уравнением (ЛОДУ)}:
    \begin{equation} \label{eq:l6_lode}
        a_0(x) y' + a_1(x) y = 0. \tag{4}
    \end{equation}
\end{definition}

\subsubsection{Алгебраические свойства и интегрирование ЛОДУ}

\begin{properties} ~
    \begin{enumerate}
        \item Если $y_1(x)$ — решение ЛОДУ \eqref{eq:l6_lode}, то для любой
        константы $\alpha \in \R$ функция $y(x) = \alpha y_1(x)$ также
        является решением ЛОДУ:
        $a_0 (\alpha y_1)' + a_1 (\alpha y_1) = \alpha (a_0 y_1' + a_1 y_1) = 0$.
        \item Если $y_1(x)$ и $y_2(x)$ — решения ЛОДУ \eqref{eq:l6_lode}, то
        их сумма $y(x) = y_1(x) + y_2(x)$ также является решением ЛОДУ:
        $a_0(y_1 + y_2)' + a_1(y_1 + y_2) = (a_0 y_1' + a_1 y_1) +
        (a_0 y_2' + a_1 y_2) = 0$.
        \item Тривиальная функция $y(x) \equiv 0$ всегда является решением
        ЛОДУ \eqref{eq:l6_lode}.
    \end{enumerate}
\end{properties}

\begin{corollary}
    Множество всех решений линейного однородного уравнения \eqref{eq:l6_lode}
    на интервале $I$ образует \textdef{линейное (векторное) пространство}.
\end{corollary}

Разделяя переменные в уравнении \eqref{eq:l6_lode} и интегрируя, получаем
\textdef{общее решение однородного уравнения} $y_{\text{оо}}(x)$:
\begin{equation} \label{eq:l6_lode_sol}
    \frac{dy}{y} = -\frac{a_1(x)}{a_0(x)} \, dx \implies
    y_{\text{оо}}(x) = C \cdot e^{-\int \frac{a_1(x)}{a_0(x)} \, dx} =
    C \cdot y_{\text{чо}}(x), \tag{5}
\end{equation}
где $y_{\text{чо}}(x) = e^{-\int \frac{a_1(x)}{a_0(x)} \, dx}$ — нетривиальное
базисное \textdef{частное решение однородного уравнения}.

\subsection{Интегрирование ЛНДУ (два базовых метода)}

\subsubsection{Метод I. Метод Лагранжа (вариация произвольной постоянной)}

\begin{proof}[Вывод формулы методом Лагранжа]~\\
    \begin{quote}
        Будем искать общее решение неоднородного уравнения \eqref{eq:l6_lnde}
        в том же виде, что и решение однородного уравнения
        \eqref{eq:l6_lode_sol}, заменив постоянную $C$ на неизвестную
        дифференцируемую функцию $C(x) \in C^1(I)$:
        \begin{equation} \label{eq:l6_lagr_sub}
            y_{\text{он}}(x) = C(x) \cdot e^{-\int \frac{a_1(x)}{a_0(x)} \, dx}.
            \tag{6}
        \end{equation}
        Вычислим производную $y_{\text{он}}'$ по правилу Лейбница и подставим
        $y_{\text{он}}$ и $y_{\text{он}}'$ в уравнение \eqref{eq:l6_lnde}:
        \begin{align*}
            y_{\text{он}} &= C(x) \cdot e^{-\int \frac{a_1(x)}{a_0(x)} \, dx}
            &&\Big| \cdot a_1(x), \\
            y_{\text{он}}' &= C'(x) \cdot e^{-\int \frac{a_1(x)}{a_0(x)} \, dx}
            - C(x) \frac{a_1(x)}{a_0(x)} e^{-\int \frac{a_1(x)}{a_0(x)} \, dx}
            &&\Big| \cdot a_0(x).
        \end{align*}
        При сложении слагаемые, содержащие $C(x)$ без производной, взаимно
        уничтожаются, и остаётся:
        \[
            a_0(x) C'(x) e^{-\int \frac{a_1(x)}{a_0(x)} \, dx} = b(x)
            \implies C(x) = \int \frac{b(x)}{a_0(x)}
            e^{\int \frac{a_1(x)}{a_0(x)} \, dx} \, dx + C.
        \]
        Подставляя найденную функцию $C(x)$ в \eqref{eq:l6_lagr_sub}, получаем
        структуру общего решения ЛНДУ:
        \begin{equation} \label{eq:l6_lnde_struct}
            y_{\text{он}}(x) =
            \underbrace{C \cdot e^{-\int \frac{a_1(x)}{a_0(x)} \, dx}}_{y_{\text{оо}}(x)}
            + \underbrace{e^{-\int \frac{a_1(x)}{a_0(x)} \, dx}
            \int \frac{b(x)}{a_0(x)} e^{\int \frac{a_1(x)}{a_0(x)} \, dx} \, dx}_{y_{\text{чн}}(x)}.
            \tag{7}
        \end{equation}
    \end{quote}
\end{proof}

\begin{theorem}[\textbf{О структуре общего решения ЛНДУ}]
    Общее решение линейного неоднородного дифференциального уравнения
    \eqref{eq:l6_lnde} равно сумме общего решения соответствующего однородного
    уравнения $y_{\text{оо}}(x)$ и любого частного решения неоднородного
    уравнения $y_{\text{чн}}(x)$:
    \[
        y_{\text{он}}(x) = y_{\text{оо}}(x) + y_{\text{чн}}(x) =
        C y_{\text{чо}}(x) + y_{\text{чн}}(x).
    \]
\end{theorem}

\begin{figure}[H]
    \centering
    \begin{tikzpicture}[scale=1.15]
        % Оси координат
        \draw[->] (-0.8, 0) -- (3.8, 0) node[right] {$x$};
        \draw[->] (0, -0.8) -- (0, 3.2) node[above] {$y$};

        % Частное решение y_чн(x) (опорная кривая)
        \draw[red, very thick, domain=-0.4:3.4, samples=80]
            plot (\x, {1.2 + 0.35*sin(120*\x)})
            node[right] {$y_{\text{чн}}(x)$};

        % Семейство решений y_он = y_чн + C * y_чо
        \foreach \c in {-1.2, -0.6, 0.6, 1.3} {
            \draw[blue, thick, dashed, domain=-0.4:3.4, samples=80]
                plot (\x, {1.2 + 0.35*sin(120*\x) + \c*exp(-0.6*\x)});
        }

        % Иллюстрация вертикального сдвига C * y_чо(x_1)
        \def\xA{0.8}
        \pgfmathsetmacro{\ybase}{1.2 + 0.35*sin(120*\xA)}
        \pgfmathsetmacro{\ytop}{\ybase + 1.3*exp(-0.6*\xA)}
        \draw[<->, teal, thick] (\xA, \ybase) -- (\xA, \ytop)
            node[midway, right] {\footnotesize $C y_{\text{чо}}(x)$};
        \fill[red] (\xA, \ybase) circle (1.5pt);
        \fill[blue] (\xA, \ytop) circle (1.5pt);
    \end{tikzpicture}
    \caption{Геометрический смысл структуры общего решения ЛНДУ: каждая
    интегральная кривая $y_{\text{он}}(x)$ получается из опорного частного
    решения $y_{\text{чн}}(x)$ прибавлением масштабированного решения
    однородного уравнения $C y_{\text{чо}}(x)$.}
\end{figure}

\subsubsection{Метод II. Метод Бернулли (представление в виде произведения)}

\begin{proof}[Вывод формулы методом Бернулли]~\\
    \begin{quote}
        Будем искать решение уравнения \eqref{eq:l6_lnde} в виде произведения
        двух неизвестных функций $u, v \in C^1(I)$:
        \begin{equation} \label{eq:l6_bern_sub}
            y(x) = u(x) \cdot v(x) \implies y' = u' v + u v'. \tag{8}
        \end{equation}
        Подставим \eqref{eq:l6_bern_sub} в уравнение \eqref{eq:l6_lnde} и
        вынесем $u(x)$ за скобки:
        \[
            a_0(x) u' v + u \cdot \underbrace{\left[a_0(x) v' +
            a_1(x) v\right]}_{= 0} = b(x).
        \]
        Выберем нетривиальную функцию $v(x)$ так, чтобы квадратная скобка
        обратилась в нуль (приняв константу интегрирования равной $1$):
        \[
            a_0(x) v' + a_1(x) v = 0 \implies
            v(x) = e^{-\int \frac{a_1(x)}{a_0(x)} \, dx}.
        \]
        Тогда для второй функции $u(x)$ остаётся уравнение:
        \[
            a_0(x) u' v(x) = b(x) \implies
            u(x) = \int \frac{b(x)}{a_0(x) v(x)} \, dx + C.
        \]
        Перемножая $y(x) = u(x) \cdot v(x)$, приходим в точности к формуле
        \eqref{eq:l6_lnde_struct}.
    \end{quote}
\end{proof}

\subsection{Уравнение Бернулли}

\begin{definition}
    \textdef{Уравнением Бернулли} называется нелинейное дифференциальное
    уравнение вида:
    \begin{equation} \label{eq:l6_bern_eq}
        a_0(x) y' + a_1(x) y = b(x) y^k, \tag{9}
    \end{equation}
    где $k \in \R \setminus \{0, 1\}$, а функции коэффициентов непрерывны на
    интервале $I \subset \R$: $a_0, a_1, b \in C(I)$, причём $a_0(x) \neq 0$
    на $I$.
\end{definition}

\begin{remark}
    При $k > 0$ тривиальная функция $y(x) \equiv 0$ всегда является решением
    уравнения \eqref{eq:l6_bern_eq} (причём при $0 < k < 1$ в силу нарушения
    условия Липшица в нуле оно является особым). При делении на $y^k$ это
    решение необходимо фиксировать отдельно.
\end{remark}

\subsubsection{Способ I. Сведение к ЛНДУ заменой переменной}

\begin{proof}[Алгоритм подстановки]~\\
    \begin{quote}
        Разделим обе части уравнения \eqref{eq:l6_bern_eq} на $y^k \neq 0$:
        \[
            a_0(x) y^{-k} y' + a_1(x) y^{1-k} = b(x).
        \]
        Введём новую неизвестную функцию:
        \begin{equation} \label{eq:l6_bern_z}
            z(x) = y^{1-k} \implies z' = (1 - k) y^{-k} y' \implies
            y^{-k} y' = \frac{z'}{1 - k}. \tag{10}
        \end{equation}
        Подстановка в разделённое уравнение приводит к ЛНДУ относительно $z(x)$:
        \begin{equation} \label{eq:l6_bern_lin}
            \frac{a_0(x)}{1 - k} z' + a_1(x) z = b(x). \tag{11}
        \end{equation}
    \end{quote}
\end{proof}

\subsubsection{Способ II. Прямое применение подстановки $y = u \cdot v$}

\begin{proof}[Алгоритм метода Бернулли]~\\
    \begin{quote}
        Подставим $y = u v$ и $y' = u' v + u v'$ непосредственно в уравнение
        \eqref{eq:l6_bern_eq}:
        \[
            a_0(x) u' v + u \cdot \underbrace{\left[a_0(x) v' +
            a_1(x) v\right]}_{= 0} = b(x) u^k v^k.
        \]
        Находя частное решение $v(x) = e^{-\int \frac{a_1(x)}{a_0(x)} \, dx}$
        из обнуления скобки, получаем для $u(x)$ уравнение с разделяющимися
        переменными:
        \[
            a_0(x) u' v(x) = b(x) u^k v(x)^k \implies
            \frac{du}{u^k} = \frac{b(x) v(x)^{k-1}}{a_0(x)} \, dx.
        \]
    \end{quote}
\end{proof}

\subsection{Уравнение Риккати и его свойства}

\begin{definition}
    \textdef{Уравнением Риккати} называется нелинейное дифференциальное
    уравнение 1-го порядка с квадратичной нелинейностью по $y$:
    \begin{equation} \label{eq:l6_riccati}
        y' = a(x) y^2 + b(x) y + c(x), \tag{12}
    \end{equation}
    где функции коэффициентов непрерывны на интервале $I \subset \R$:
    $a, b, c \in C(I)$ и $a(x) \not\equiv 0$.
\end{definition}

\begin{remark}
    В общем случае уравнение Риккати \eqref{eq:l6_riccati} не интегрируется в
    квадратурах. По теореме Лиувилля специальное уравнение Риккати
    \begin{equation} \label{eq:l6_riccati_spec}
        y' = y^2 + x^\alpha \tag{13}
    \end{equation}
    интегрируется в элементарных функциях тогда и только тогда, когда
    $\alpha = -2$ или $\alpha = \dfrac{-4k}{2k \pm 1}$, где $k \in \Z$.
\end{remark}

\subsubsection{Случай постоянных коэффициентов}

Если коэффициенты $a, b, c$ постоянны, то уравнение Риккати \eqref{eq:l6_riccati}
представляет собой автономное уравнение с разделяющимися переменными.
Характерный пример $y' = y^2 - 2y + 1 = (y - 1)^2$ был подробно решён и
проиллюстрирован в Лекции 3 (см. уравнение \eqref{eq:l3_riccati_ex} и рис. к
нему): стационарное решение $y \equiv 1$ определено на всём $\R$, тогда как
все остальные решения $y(x) = 1 - \frac{1}{x + C}$ разрушаются на вертикальных
асимптотах $x = -C$ из-за квадратичного роста правой части.

\subsubsection{Сведение уравнения Риккати к линейному}

\begin{theorem}[\textbf{Об интегрировании уравнения Риккати}]
    Если известно хотя бы одно частное решение $y_1(x) \in C^1(I)$ уравнения
    Риккати \eqref{eq:l6_riccati}, то нахождение его общего решения сводится к
    интегрированию ЛНДУ 1-го порядка.
\end{theorem}

\begin{proof}[Обоснование метода]~\\
    \begin{quote}
        Выполним замену искомой функции, отсчитывая отклонение $z(x)$ от
        известного частного решения $y_1(x)$:
        \begin{equation} \label{eq:l6_riccati_sub}
            y(x) = y_1(x) + z(x) \implies y' = y_1' + z'. \tag{14}
        \end{equation}
        Подставим \eqref{eq:l6_riccati_sub} в уравнение \eqref{eq:l6_riccati} и
        раскроем скобки:
        \[
            y_1' + z' = \left[a(x) y_1^2 + b(x) y_1 + c(x)\right] +
            \left(b(x) + 2 a(x) y_1(x)\right) z + a(x) z^2.
        \]
        Поскольку $y_1(x)$ — решение исходного уравнения, квадратная скобка
        сокращается с $y_1'$, и мы получаем \textdef{уравнение Бернулли}
        относительно $z(x)$ с показателем $k = 2$:
        \begin{equation} \label{eq:l6_riccati_bern}
            z' - \left(b(x) + 2 a(x) y_1(x)\right) z = a(x) z^2. \tag{15}
        \end{equation}
        Применяя стандартную подстановку Бернулли \eqref{eq:l6_bern_z} при
        $k = 2$:
        \begin{equation} \label{eq:l6_riccati_u}
            u(x) = z^{-1} = \frac{1}{z(x)} \implies \frac{z'}{z^2} = -u',
            \tag{16}
        \end{equation}
        получаем \textdef{линейное неоднородное уравнение (ЛНДУ)} относительно
        $u(x)$:
        \begin{equation} \label{eq:l6_riccati_lin}
            u' + \left(b(x) + 2 a(x) y_1(x)\right) u = -a(x). \tag{17}
        \end{equation}
        Следовательно, сквозная подстановка $y(x) = y_1(x) + \dfrac{1}{u(x)}$
        сводит уравнение Риккати непосредственно к ЛНДУ \eqref{eq:l6_riccati_lin}.
    \end{quote}
\end{proof}

\subsubsection{Задачи о кратных частных решениях уравнения Риккати}

\begin{problem}
    Пусть $y_1(x), y_2(x)$ — два различных частных решения уравнения Риккати
    \eqref{eq:l6_riccati}. Доказать, что общее решение уравнения находится с
    помощью \textbf{одного интегрирования}.
\end{problem}

\begin{proof}[Решение]~\\
    \begin{quote}
        Подстановка $y = y_1(x) + \frac{1}{u(x)}$ приводит уравнение Риккати к
        ЛНДУ \eqref{eq:l6_riccati_lin}:
        \begin{equation} \label{eq:l6_task1_lin}
            u' + P(x) u = Q(x), \tag{18}
        \end{equation}
        где $P(x) = b(x) + 2 a(x) y_1(x)$ и $Q(x) = -a(x)$. Так как известно
        второе решение $y_2(x) \neq y_1(x)$, функция
        \[
            u_1(x) = \frac{1}{y_2(x) - y_1(x)}
        \]
        является известным частным решением ЛНДУ \eqref{eq:l6_task1_lin}:
        \begin{equation} \label{eq:l6_task1_u1}
            u_1' + P(x) u_1 = Q(x). \tag{19}
        \end{equation}
        По теореме 6.1 о структуре общего решения ЛНДУ положим
        $u(x) = u_1(x) + v(x)$. Вычитая тождество \eqref{eq:l6_task1_u1} из
        \eqref{eq:l6_task1_lin}, получаем для $v(x)$ однородное уравнение
        (ЛОДУ):
        \begin{equation} \label{eq:l6_task1_v}
            v' + P(x) v = 0 \implies v(x) = C \cdot e^{-\int P(x) \, dx}.
            \tag{20}
        \end{equation}
        Для нахождения $v(x)$ требуется взять ровно один интеграл, после чего
        общее решение записывается как
        \[
            y(x) = y_1(x) + \frac{1}{\frac{1}{y_2(x) - y_1(x)} +
            C \cdot e^{-\int P(x) \, dx}}.
        \]
    \end{quote}
\end{proof}

\begin{problem}
    Пусть $y_1(x), y_2(x), y_3(x)$ — три попарно различных частных решения
    уравнения Риккати \eqref{eq:l6_riccati}. Доказать, что общее решение
    находится чисто алгебраически — \textbf{без единого интегрирования}.
\end{problem}

\begin{proof}[Решение]~\\
    \begin{quote}
        Для линейного уравнения \eqref{eq:l6_task1_lin} нам теперь известны два
        различных частных решения:
        \[
            u_1(x) = \frac{1}{y_2(x) - y_1(x)}, \quad
            u_2(x) = \frac{1}{y_3(x) - y_1(x)}.
        \]
        Их разность $v_0(x) = u_1(x) - u_2(x) \neq 0$ удовлетворяет
        соответствующему однородному уравнению (ЛОДУ):
        \begin{equation} \label{eq:l6_task2_hom}
            (u_1 - u_2)' + P(x) (u_1 - u_2) = 0. \tag{21}
        \end{equation}
        Следовательно, общее решение ЛОДУ имеет вид $v(x) = C \cdot (u_1 - u_2)$,
        а общее решение ЛНДУ \eqref{eq:l6_task1_lin} записывается без
        интегрирования как:
        \[
            u(x) = u_1(x) + C \cdot \left(u_1(x) - u_2(x)\right).
        \]
        Возвращаясь к $y(x) = y_1(x) + \frac{1}{u(x)}$, получаем общее решение
        уравнения Риккати без квадратур.
    \end{quote}
\end{proof}

\begin{remark}
    Подставляя выражения для $u, u_1, u_2$ в соотношение
    $\frac{u - u_1}{u_2 - u_1} = \tilde{C}$, мы приходим к теореме о
    \textdef{постоянстве сложного (ангармонического) отношения} любых четырёх
    решений уравнения Риккати:
    \[
        \frac{y - y_2}{y - y_1} : \frac{y_3 - y_2}{y_3 - y_1} = \tilde{C}.
    \]
\end{remark}